\ekddiv{type=act, n=1, depth=2}
\stage{ACTE 1}
\ekddiv{type=scene, n=1, depth=3}
\stage{SCENE III.}
\stage{IOCASTE, ANTIGONE, ETEOCLE, OLYMPE.}
\begin{speech}
  \speaker{IOCASTE.}
  \begin{ekdverse}
    \vnum{1}{OLympe, soustiens~moy, ma douleur est extréme.\\}
  \end{ekdverse}
\end{speech}
\begin{speech}
  \speaker{ETEOCLE.}
  \begin{ekdverse}
    \vnum{2}{\app{\lem[wit={A}]{Madame qu’auez~vous~? Et quel mal si caché…}\rdg[wit={B}]{Madame, qu’avez-vous~? Et quel mal si caché…}\rdg[wit={C}]{Madame, qu’avez-vous~? \& quel mal si caché…}\rdg[wit={D}]{Madame, qu’avez-vous~? \& quel trouble…}\rdg[wit={E}]{Madame qu’avez~vous~? Et quel mal si caché….}}\\}
  \end{ekdverse}
\end{speech}
\begin{speech}
  \speaker{IOCASTE.}
  \begin{ekdverse}
    \vnum{3}{\app{\lem[wit={A, E}]{Ah~! mon Fils, de quel sang estes vous-là taché~?}\rdg[wit={B, C}]{Ah~! mon Fils, de quel sang revenez-vous taché~?}\rdg[wit={D}]{~~~~~~~~~~~~~~~~~~Ah~! mon Fils,}}\\}
    \vnum{4}{\app{\lem[wit={A, B, C, D}]{(lacune)}\rdg[wit={E}]{Quelles traces de sang voi-je sur vos habits~?}}\\}
    \vnum{5}{\app{\lem[wit={A, E}]{Est-ce~de~vostre~Frere,}\rdg[wit={B, C, D}]{Est-ce~du~sang~d’un~Frere,}} ou n’est-ce point du vostre~?\\}
  \end{ekdverse}
\end{speech}
\begin{speech}
  \speaker{ETEOCLE.}
  \begin{ekdverse}
    \vnum{6}{Non, Madame, ce n’est ny de l’vn ny de l’autre,\\}
    \vnum{7}{\app{\lem[wit={A, B, C, E}]{Polinice à mes yeux ne s’est point presenté,}\rdg[wit={D}]{Dans son camp jusqu’icy Polinice arresté,}}\\}
    \vnum{8}{\app{\lem[wit={A}]{Et l’on s’est peu battu d’vn \& d’autre costé.}\rdg[wit={B, C}]{Et l’on s’est peu battu d’un \& d’autre costé,}\rdg[wit={D}]{Pour combattre à mes yeux ne s’est point presenté.}\rdg[wit={E}]{Et l’on s’est peu battu d’un \& d’autre costé.}}\\}
    \vnum{9}{\app{\lem[wit={A}]{Seulement quelques Grecs d’vn insolent courage,}\rdg[wit={B}]{Du Camp des Argiens une troupe hardie,}\rdg[wit={C}]{Du Camp des Argiens une troupe hardie}\rdg[wit={D}]{D’Argiens seulement une troupe hardie}\rdg[wit={E}]{Seulement quelques Grecs d’un insolent courage,}}\\}
    \vnum{10}{\app{\lem[wit={A, E}]{M’ayant osé d’abord disputer le passage,}\rdg[wit={B, C, D}]{M’a voulu de nos murs disputer la sortie.}}\\}
    \vnum{11}{I’ay fait mordre la poudre à ces audacieux,\\}
    \vnum{12}{Et leur sang est celuy qui paroist à vos yeux.\\}
  \end{ekdverse}
\end{speech}
\begin{speech}
  \speaker{IOCASTE.}
  \begin{ekdverse}
    \vnum{13}{\app{\lem[wit={A}]{Mais pourquoy donc sortir auecque vostre Armée,}\rdg[wit={B}]{Mais pourquoy donc sortir auecque vostre armée~?}\rdg[wit={C}]{Mais pourquoi donc sortir avecque vostre armée~?}\rdg[wit={D}]{Mais que prétendiez vous~? \& quelle ardeur soudaine}\rdg[wit={E}]{Mais pourquoy donc sortir avecque vostre Armée,}}\\}
    \vnum{14}{\app{\lem[wit={A}]{Quel est ce mouuement qui m’a tant alarmée~?}\rdg[wit={B, C, E}]{Quel est ce mouvement qui m’a tant alarmée~?}\rdg[wit={D}]{Vous a fait tout à coup descendre dans la plaine~?}}\\}
  \end{ekdverse}
\end{speech}
\begin{speech}
  \speaker{ETEOCLE.}
  \begin{ekdverse}
    \vnum{15}{Madame il estoit temps que j’en vsasse ainsi,\\}
    \vnum{16}{Et ie perdois ma gloire à demeurer icy.\\}
    \vnum{17}{\app{\lem[wit={A}]{Ie n’ay que trop languy derriere vne muraille,}\rdg[wit={B, C}]{Je n’ay que trop langui derriere une muraille,}\rdg[wit={D}]{(lacune)}\rdg[wit={E}]{Je n’ay que trop languy derriere une muraille,}}\\}
    \vnum{18}{\app{\lem[wit={A}]{Ie brûlois de me voir en vn champ de bataille,}\rdg[wit={B, C}]{Je brûlois de me voir en un champ de bataille.}\rdg[wit={D}]{(lacune)}\rdg[wit={E}]{Je brûlois de me voir en un champ de bataille,}}\\}
    \vnum{19}{\app{\lem[wit={A, B, C, E}]{Lors~que l’on peut paroistre au milieu des hazards,}\rdg[wit={D}]{(lacune)}}\\}
    \vnum{20}{\app{\lem[wit={A}]{Vn grand cœur est honteux de garder des remparts.}\rdg[wit={B, C}]{Un grand cœur est honteux de garder des rempars.}\rdg[wit={D}]{(lacune)}\rdg[wit={E}]{Un grand cœur est honteux de garder des remparts.}}\\}
    \vnum{21}{\app{\lem[wit={A}]{I’estois las d’endurer que le fier Polinice}\rdg[wit={B, C, E}]{J’estois las d’endurer que le fier Polinice}\rdg[wit={D}]{(lacune)}}\\}
    \vnum{22}{\app{\lem[wit={A, B, C, E}]{Me reprochast tout haut cét indigne exercice,}\rdg[wit={D}]{(lacune)}}\\}
    \vnum{23}{\app{\lem[wit={A, E}]{Et criast aux Thebains, afin de les gaigner,}\rdg[wit={B, C}]{Et criast aux Thebains, afin de les gagner,}\rdg[wit={D}]{(lacune)}}\\}
    \vnum{24}{\app{\lem[wit={A, B}]{Que ie laissois aux fers ceux qui me font regner.}\rdg[wit={C}]{Que je laissois perir ceux qui me font regner.}\rdg[wit={D}]{(lacune)}\rdg[wit={E}]{Que je laissois aux fers ceux qui me font regner.}}\\}
    \vnum{25}{Le Peuple à qui la faim se faisoit déja craindre,\\}
    \vnum{26}{De mon peu de vigueur commençoit à se plaindre,\\}
    \vnum{27}{Me reprochant déja qu’il m’auoit couronné,\\}
    \vnum{28}{Et que j’occupois mal le rang qu’il m’a donné.\\}
    \vnum{29}{Il le faut satisfaire, \& quoy~qu’il en arriue,\\}
    \vnum{30}{Thebes dés aujourd’huy ne sera plus captiue,\\}
    \vnum{31}{Ie veux, en n’y laissant aucun de mes soldats,\\}
    \vnum{32}{Qu’elle soit seulement juge de nos combats.\\}
    \vnum{33}{I’ay des forces assez pour tenir la \app{\lem[wit={A, C, D, E}]{campagne,}\rdg[wit={B}]{camgagne,}}\\}
    \vnum{34}{Et si quelque bon-heur nos armes accompagne,\\}
    \vnum{35}{L’insolent Polinice \& ses \app{\lem[wit={A, E}]{Grecs}\rdg[wit={B, C, D}]{fiers}} \app{\lem[wit={A, E}]{orgüeilleux,}\rdg[wit={B}]{Alliez,}\rdg[wit={C, D}]{Alliez}}\\}
    \vnum{36}{Laisseront Thebes libre, ou mouront à \app{\lem[wit={A, E}]{ses}\rdg[wit={B, C, D}]{mes}} \app{\lem[wit={A, E}]{yeux.}\rdg[wit={B, C, D}]{piez.}}\\}
  \end{ekdverse}
\end{speech}
\begin{speech}
  \speaker{IOCASTE.}
  \begin{ekdverse}
    \vnum{37}{\app{\lem[wit={A}]{Vous preserue le Ciel d’vne telle Victoire,}\rdg[wit={B}]{Vous preserve le Ciel d’une telle Victoire.}\rdg[wit={C}]{Vous preserve le Ciel d’une telle victoire.}\rdg[wit={D}]{(lacune)}\rdg[wit={E}]{Vous preserve le Ciel d’une telle Victoire,}}\\}
    \vnum{38}{\app{\lem[wit={A}]{Thebes ne veut point voir vne action si noire,}\rdg[wit={B, C, E}]{Thebes ne veut point voir une action si noire,}\rdg[wit={D}]{(lacune)}}\\}
    \vnum{39}{\app{\lem[wit={A, E}]{Laissez là son salut \& n’y songez jamais~;}\rdg[wit={B, C}]{Laissez-là son salut \& n’y songez jamais~;}\rdg[wit={D}]{(lacune)}}\\}
    \vnum{40}{\app{\lem[wit={A, B, E}]{La Guerre vaut bien mieux que cette affreuse Paix.}\rdg[wit={C}]{La Guerre vaut bien-mieux que cette affreuse Paix.}\rdg[wit={D}]{(lacune)}}\\}
    \vnum{41}{\app{\lem[wit={A}]{Dure-t’elle à iamais cette cruelle Guerre,}\rdg[wit={B}]{Dure-t’elle à jamais cette cruelle Guerre,}\rdg[wit={C}]{(lacune)}\rdg[wit={D, E}]{Dont le flambeau fatal desole cette terre.}}\\}
    \vnum{42}{\app{\lem[wit={A}]{Prolongez nos mal-heurs, augmentez-les tousiours,}\rdg[wit={B}]{Prolongez nos mal-heurs, augmentez-les tousjours,}\rdg[wit={C}]{(lacune)}\rdg[wit={D}]{Plustost qu’vn si grand crime en arreste le cours.}\rdg[wit={E}]{Plustost qu’un si grand crime en arreste le cours.}}\\}
    \vnum{43}{\app{\lem[wit={A}]{Vous-mesme d’vn tel sang soüilleriez vous vos Armes~?}\rdg[wit={B}]{Vous-mesme d’un tel sang soüilleriez-vous vos Armes~?}\rdg[wit={C}]{Vous-mesme d’un tel sang, soüillerez-vous vos Armes~?}\rdg[wit={D}]{Vous pourriez d’un tel sang, ô Ciel~! souiller vos armes~?}\rdg[wit={E}]{Vous-mesme d’un tel sang soüilleriez vous vos Armes~?}}\\}
    \vnum{44}{La Courone pour vous a~t’elle tant de charmes~?\\}
    \vnum{45}{Si par vn parricide il la falloit gaigner\\}
    \vnum{46}{Ah~! mon Fils, à ce prix voudriez~vous regner~?\\}
    \vnum{47}{Mais il ne tient qu’à vous si l’honneur vous anime,\\}
    \vnum{48}{De nous donner la Paix, sans le secours d’vn crime,\\}
    \vnum{49}{\app{\lem[wit={A}]{Vous pouuez-vous montrer genereux tout a fait,}\rdg[wit={B}]{Vous pouvez-vous montrer genereux tout à fait,}\rdg[wit={C}]{Vous pouvez vous montrer genereux tout à fait,}\rdg[wit={D}]{Et de vostre couroux triomphant aujourd’huy}\rdg[wit={E}]{Vous pouvez-vous montrer genereux tout a fait,}}\\}
    \vnum{50}{Contenter vostre Frere, \& regner \app{\lem[wit={A, B, C, E}]{en}\rdg[wit={D}]{avec}} \app{\lem[wit={A, B, C, E}]{effet.}\rdg[wit={D}]{luy.}}\\}
  \end{ekdverse}
\end{speech}
\begin{speech}
  \speaker{ETEOCLE.}
  \begin{ekdverse}
    \vnum{51}{Appellez-vous regner \app{\lem[wit={A, E}]{luy~ceder}\rdg[wit={B, C}]{lui~ceder}\rdg[wit={D}]{partager}} ma Couronne,\\}
    \vnum{52}{\app{\lem[wit={A, B, C, E}]{Quand le sang, \& le Peuple à la fois me la donne~?}\rdg[wit={D}]{Et ceder laschement ce que mon droit me donne~?}}\\}
  \end{ekdverse}
\end{speech}
\begin{speech}
  \speaker{IOCASTE.}
  \begin{ekdverse}
    \vnum{53}{\app{\lem[wit={A}]{Vous sçauez bien, mon Fils, que le choix \& le sang}\rdg[wit={B, C, D}]{Vous le sçavez, mon Fils, la justice \& le sang}\rdg[wit={E}]{Vous sçavez bien, mon Fils, que le choix \& le sang}}\\}
    \vnum{54}{Luy \app{\lem[wit={A, B, D, E}]{donnent}\rdg[wit={C}]{donne}} comme à vous sa part à ce haut rang.\\}
    \vnum{55}{Oedipe en acheuant sa triste destinée\\}
    \vnum{56}{Ordonna que chacun regneroit son année,\\}
    \vnum{57}{Et n’ayant qu’vn Estat à mettre sous vos Lois,\\}
    \vnum{58}{\app{\lem[wit={A, B, C, E}]{Il voulut que tous deux vous en fussiez les Rois.}\rdg[wit={D}]{Voulut que tour à tour vous fussiez tous deux Rois.}}\\}
    \vnum{59}{A ces conditions vous \app{\lem[wit={A, B, C, E}]{voulustes}\rdg[wit={D}]{daignastes}} souscrire,\\}
    \vnum{60}{Le sort vous appella le premier à l’Empire,\\}
    \vnum{61}{Vous montastes au Trosne, il n’en fut point jaloux,\\}
    \vnum{62}{Et vous ne voulez pas qu’il y monte apres vous~?\\}
  \end{ekdverse}
\end{speech}
\begin{speech}
  \speaker{\app{\lem[wit={A, E}]{ETEOCLE.}\rdg[wit={B, C, D}]{(lacune)}}}
  \begin{ekdverse}
    \vnum{63}{\app{\lem[wit={A, E}]{Il est vray, je promis, ce que voulut mon Pere,}\rdg[wit={B, C, D}]{(lacune)}}\\}
    \vnum{64}{\app{\lem[wit={A}]{Pour vn Trosne est-il rien qu’on refuse de faire~?}\rdg[wit={B, C, D}]{(lacune)}\rdg[wit={E}]{Pour un Trosne est-il rien qu’on refuse de faire~?}}\\}
    \vnum{65}{\app{\lem[wit={A}]{On promet tout, Madame, afin d’y paruenir,}\rdg[wit={B, C, D}]{(lacune)}\rdg[wit={E}]{On promet tout, Madame, afin d’y parvenir,}}\\}
    \vnum{66}{\app{\lem[wit={A, E}]{Mais on ne songe apres qu’à s’y bien maintenir.}\rdg[wit={B, C, D}]{(lacune)}}\\}
    \vnum{67}{\app{\lem[wit={A}]{I’estois alors sujet, \& dans l’obeïssance,}\rdg[wit={B, C, D}]{(lacune)}\rdg[wit={E}]{J’estois alors sujet, \& dans l’obeïssance,}}\\}
    \vnum{68}{\app{\lem[wit={A, E}]{Et je tiens aujourd’huy la supréme puissance :}\rdg[wit={B, C, D}]{(lacune)}}\\}
    \vnum{69}{\app{\lem[wit={A}]{Ce que ie fis alors ne m’est plus vne Loy,}\rdg[wit={B, C, D}]{(lacune)}\rdg[wit={E}]{Ce que je fis alors ne m’est plus une Loy,}}\\}
    \vnum{70}{\app{\lem[wit={A}]{Le deuoir d’vn Sujet n’est pas celuy d’vn Roy.}\rdg[wit={B, C, D}]{(lacune)}\rdg[wit={E}]{Le devoir d’un Sujet n’est pas celuy d’un Roy.}}\\}
    \vnum{71}{\app{\lem[wit={A, E}]{D’abord que sur sa teste il reçoit la Couronne,}\rdg[wit={B, C, D}]{(lacune)}}\\}
    \vnum{72}{\app{\lem[wit={A}]{Vn Roy sort à l’instant de sa propre personne,}\rdg[wit={B, C, D}]{(lacune)}\rdg[wit={E}]{Un Roy sort à l’instant de sa propre personne,}}\\}
    \vnum{73}{\app{\lem[wit={A}]{L’interest du public doit deuenir le sien,}\rdg[wit={B, C, D}]{(lacune)}\rdg[wit={E}]{L’interest du public doit devenir le sien,}}\\}
    \vnum{74}{\app{\lem[wit={A, E}]{Il doit tout à l’Estat, \& ne se doit plus rien.}\rdg[wit={B, C, D}]{(lacune)}}\\}
  \end{ekdverse}
\end{speech}
\begin{speech}
  \speaker{\app{\lem[wit={A}]{IOCASTE.}\rdg[wit={B, C, D}]{(lacune)}\rdg[wit={E}]{JOCASTE.}}}
  \begin{ekdverse}
    \vnum{75}{\app{\lem[wit={A, E}]{Au moins doit-il, mon Fils, quelque chose à sa gloire,}\rdg[wit={B, C, D}]{(lacune)}}\\}
    \vnum{76}{\app{\lem[wit={A, E}]{Dont le soin ne doit pas sortir de sa memoire,}\rdg[wit={B, C, D}]{(lacune)}}\\}
    \vnum{77}{\app{\lem[wit={A}]{Et quand ce nouueau rang l’affranchiroit des Lois,}\rdg[wit={B, C, D}]{(lacune)}\rdg[wit={E}]{Et quand ce nouveau rang l’affranchiroit des Lois,}}\\}
    \vnum{78}{\app{\lem[wit={A, E}]{Au moins doit il tenir sa parole à des Rois.}\rdg[wit={B, C, D}]{(lacune)}}\\}
  \end{ekdverse}
\end{speech}
\begin{speech}
  \speaker{ETEOCLE.}
  \begin{ekdverse}
    \vnum{79}{\app{\lem[wit={A, E}]{Polinice à ce titre auroit tort de pretendre,}\rdg[wit={B, C}]{Non, Madame, à l’Empire il ne doit plus prétendre.}\rdg[wit={D}]{Non, Madame, à l’Empire il ne doit plus prétendre~:}}\\}
    \vnum{80}{\app{\lem[wit={A}]{Thebes~sous~son~pouuoir}\rdg[wit={B, C, E}]{Thebes~sous~son~pouvoir}\rdg[wit={D}]{Thebes~à~cet~arrest}} n’a point voulu se rendre,\\}
    \vnum{81}{Et lors~que sur le Trosne il s’est voulu placer,\\}
    \vnum{82}{C’est elle \& non pas moy qui l’en a sçeû chasser.\\}
    \vnum{83}{Thebes doit~elle moins redouter sa puissance,\\}
    \vnum{84}{Apres auoir six mois senty sa violence~?\\}
    \vnum{85}{Voudroit-elle obeïr à ce Prince inhumain,\\}
    \vnum{86}{Qui vient d’armer \app{\lem[wit={A, E}]{contr’-elle}\rdg[wit={B, C, D}]{contre~elle}} \& le fer \& la faim~?\\}
    \vnum{87}{Prendroit-elle pour Roy l’Esclaue de Mycéne,\\}
    \vnum{88}{Qui pour tous les Thebains n’a plus que de la haine,\\}
    \vnum{89}{Qui s’est au Roy d’Argos indignement soûmis,\\}
    \vnum{90}{Et que l’Hymen attache à nos fiers ennemis~?\\}
    \vnum{91}{Lors~que le Roy d’Argos l’a choisi pour son Gendre,\\}
    \vnum{92}{Il esperoit par luy de voir Thebes en cendre,\\}
    \vnum{93}{L’amour eût peu de part à cét hymen honteux,\\}
    \vnum{94}{Et la seule fureur en alluma les feux.\\}
    \vnum{95}{Thebes m’a couronné pour éuiter ses \app{\lem[wit={A, E}]{chaisnes~;}\rdg[wit={B}]{chaisnes.}\rdg[wit={C}]{chaisnes,}\rdg[wit={D}]{chaînes,}}\\}
    \vnum{96}{Elle s’attend par moy de voir finir ses peines,\\}
    \vnum{97}{Il la faut accuser si je manque de Foy,\\}
    \vnum{98}{Et je suis son Captif, ie ne suis pas son Roy.\\}
  \end{ekdverse}
\end{speech}
\begin{speech}
  \speaker{IOCASTE.}
  \begin{ekdverse}
    \vnum{99}{Dites, dites plutost, cœur ingrat \& farouche,\\}
    \vnum{100}{Qu’auprés du Diadême il n’est rien qui vous touche~;\\}
    \vnum{101}{Mais ie me trompe encor ce rang ne vous plaist pas,\\}
    \vnum{102}{Et le crime tout seul a pour vous des appas.\\}
    \vnum{103}{Hé~!~bien, puis~qu’à ce point vous en estes auide,\\}
    \vnum{104}{Ie vous offre à commettre vn double parricide,\\}
    \vnum{105}{Versez le sang d’vn Frere~: Et si c’est peu du sien,\\}
    \vnum{106}{Ie vous inuite encore à \app{\lem[wit={A, E}]{respandre}\rdg[wit={B, C, D}]{répandre}} le mien.\\}
    \vnum{107}{Vous n’aurez plus alors d’ennemis à \app{\lem[wit={A, E}]{sousmettre,}\rdg[wit={B, C}]{soûmettre,}\rdg[wit={D}]{soumettre,}}\\}
    \vnum{108}{D’obstacle à surmonter ny de crime à commettre,\\}
    \vnum{109}{Et n’ayant plus au Trosne vn \app{\lem[wit={A, E}]{fascheux}\rdg[wit={B, C, D}]{fâcheux}} \app{\lem[wit={A, D, E}]{concurrent,}\rdg[wit={B, C}]{concurrant,}}\\}
    \vnum{110}{De tous les Criminels vous serez le plus grand,\\}
  \end{ekdverse}
\end{speech}
\begin{speech}
  \speaker{ETEOCLE.}
  \begin{ekdverse}
    \vnum{111}{Hé bien, Madame, hé bien, il faut vous satisfaire,\\}
    \vnum{112}{Il faut sortir du Trosne \& couronner mon frere,\\}
    \vnum{113}{Il faut pour seconder vostre injuste projet,\\}
    \vnum{114}{De son Roy que j’estois deuenir son sujet.\\}
    \vnum{115}{Et pour vous éleuer au comble de la joye,\\}
    \vnum{116}{Il faut à sa fureur que ie me liure en proye,\\}
    \vnum{117.1}{Il faut par mon trépas….\\}
  \end{ekdverse}
\end{speech}
\begin{speech}
  \speaker{IOCASTE.}
  \begin{ekdverse}
    \vnum{117.2}{Ah Ciel~! quelle rigueur,\\}
    \vnum{118}{Que vous penetrez mal dans le \app{\lem[wit={A, B, C, E}]{fonds}\rdg[wit={D}]{fond}} de mon cœur~!\\}
    \vnum{119}{Ie ne demande pas que vous quittiez l’Empire,\\}
    \vnum{120}{Regnez \app{\lem[wit={A}]{tousiours,}\rdg[wit={B, C, D}]{toûjours,}\rdg[wit={E}]{tousjours,}} mon Fils, c’est ce que ie desire.\\}
    \vnum{121}{Mais si tant de malheurs vous touchent de pitié,\\}
    \vnum{122}{Si pour moy vostre cœur garde quelque amitié~;\\}
    \vnum{123}{Et si vous prenez soin de vostre gloire mesme,\\}
    \vnum{124}{Associez un Frere à cét honneur supréme~;\\}
    \vnum{125}{Ce n’est qu’vn vain éclat qu’il receura de vous,\\}
    \vnum{126}{Vostre regne en sera plus puissant \& plus doux.\\}
    \vnum{127}{Les Peuples admirant cette vertu sublime,\\}
    \vnum{128}{Voudront \app{\lem[wit={A}]{tousiours}\rdg[wit={B, C, D}]{toûjours}\rdg[wit={E}]{tousjours}} pour Prince vn Roy si magnanime,\\}
    \vnum{129}{Et cét illustre effort, loin d’affoiblir vos droits,\\}
    \vnum{130}{Vous rendra le plus juste \& le plus grand des Rois.\\}
    \vnum{131}{Ou s’il faut que mes vœux vous \app{\lem[wit={A}]{treuuent}\rdg[wit={B, C, D}]{trouvent}\rdg[wit={E}]{treuvent}} inflexible,\\}
    \vnum{132}{Si la Paix à ce prix vous paroist impossible,\\}
    \vnum{133}{Et \app{\lem[wit={A, E}]{que}\rdg[wit={B, C, D}]{si}} le Diadéme \app{\lem[wit={A, E}]{ait}\rdg[wit={B, C, D}]{a}} pour vous tant d’attraits,\\}
    \vnum{134}{Au moins consolez-moy de quelque heure de Paix,\\}
    \vnum{135}{\app{\lem[wit={A, B, C, E}]{Accordez quelque tréve à ma douleur amere,}\rdg[wit={D}]{Accordez cette grace aux larmes d’une Mere.}}\\}
    \vnum{136}{Et cependant, mon Fils, j’iray voir vostre Frere,\\}
    \vnum{137}{La pitié dans son ame aura peut~estre lieu,\\}
    \vnum{138}{Ou du moins pour jamais j’iray luy dire Adieu.\\}
    \vnum{139}{Dés ce mesme moment permettez que ie sorte,\\}
    \vnum{140}{I’iray jusqu’à sa tente, \& j’iray sans escorte,\\}
    \vnum{141}{\app{\lem[wit={A}]{Dans cette occasion rien ne peut m’émouuoir.}\rdg[wit={B, C, E}]{Dans cette occasion rien ne peut m’émouvoir.}\rdg[wit={D}]{Par mes justes soupirs j’espere l’émouvoir.}}\\}
  \end{ekdverse}
\end{speech}
\begin{speech}
  \speaker{ETEOCLE.}
  \begin{ekdverse}
    \vnum{142}{Madame, sans sortir vous le pouuez \app{\lem[wit={A, E}]{bien~voir.}\rdg[wit={B, C}]{revoir.}\rdg[wit={D}]{revoir~;}}\\}
    \vnum{143}{Et si cette entreueuë à pour vous tant de charmes,\\}
    \vnum{144}{Il ne tiendra qu’à luy de suspendre nos armes,\\}
    \vnum{145}{Vous pouuez dés cette heure accomplir vos souhaits,\\}
    \vnum{146}{Et le faire venir jusques dans ce Palais.\\}
    \vnum{147}{\app{\lem[wit={A}]{Ie~feray~plus~encore,}\rdg[wit={B}]{Je~feray~plus~encor,}\rdg[wit={C}]{Je~ferai~plus~encore,}\rdg[wit={D}]{J’iray~plus~loin~encore,}\rdg[wit={E}]{Je~feray~plus~encore,}} \& pour faire \app{\lem[wit={A, E}]{connoistre,}\rdg[wit={B, C, D}]{connaistre,}}\\}
    \vnum{148}{Qu’il a tort en effet de me nommer vn traistre,\\}
    \vnum{149}{Et que ie ne suis pas vn tyran odieux,\\}
    \vnum{150}{Que l’on fasse parler \& le Peuple \& les Dieux.\\}
    \vnum{151}{Si le Peuple \app{\lem[wit={A, B, C, E}]{le~veut,}\rdg[wit={D}]{y~consent,}} ie luy cede ma place,\\}
    \vnum{152}{Mais qu’il se rende \app{\lem[wit={A, E}]{aussi}\rdg[wit={B}]{aussy}\rdg[wit={C, D}]{enfin}} si le Peuple le chasse,\\}
    \vnum{153}{Ie ne force personne, \& j’engage ma foy\\}
    \vnum{154}{De laisser aux Thebains à se choisir vn Roy.\\}
  \end{ekdverse}
\end{speech}
